\properlane{one-people}{The one Lord gathers a people who bear one another}

The Mass gives the assembly more than a collection of individual religious tasks. The Collect addresses the needs of God's people, Paul names one body, the Gradual sings election, Daniel intercedes for a sanctuary and a people, and the final prayers ask for healing in the plural. Even the singular voice of the Introit stands beside the blessed walkers of God's law. A person approaches God as a servant whose need is personal, yet whose good belongs within a common calling.

That calling confronts a familiar difficulty. People may prefer a community made only of those whose habits, abilities, and opinions are convenient for them. Paul's appeal begins elsewhere. The one body already owes its life to one Spirit and has been called to one hope. Patience is necessary because the members are not chosen by each other's convenience. The common gift becomes the ground of the demanding work of bearing with one another.

\subsection{Unity received before it is maintained}
Chrysostom's exposition of Ephesians 4:3 gives the Spirit an active role in joining those divided by age, wealth, sex, and manner of life. Spiritual kinship is stronger than those differences. His image of fire joining separate pieces into one blaze gives unity warmth and movement: charity binds because persons are being drawn into a shared life. But discord can obstruct that life. The diligence Paul commands is not needless effort after a gift has been given; it is the care appropriate to a gift that persons can injure.\footnote{Chrysostom, \textit{Ephesians}, hom.~9, exposition of Eph.~4:3.}

The next homily widens the body across place and time. The faithful now living, those who have lived, and those who will live belong to the one body; Chrysostom includes those who pleased God before Christ's coming through their relation to him. The congregation is therefore not the whole measure of the Church. Its present quarrels are set within a communion that exceeds its local memory. No member can make the body equivalent to the people he happens to recognize.\footnote{Chrysostom, \textit{Ephesians}, hom.~10, opening exposition of Eph.~4:4--5.}

Chrysostom also offers several ways to understand the one spirit, including the received Spirit and the concord or zeal of those joined together. His explanations are related, not interchangeable definitions. In his account the divine giver and the common disposition of his people are distinct, so a group's intense shared enthusiasm does not by itself prove that its purpose is holy.

Thomas makes the difference explicit in his account of spiritual unity. People can unite in evil; the unity sought by the apostle is ordered to a good end. Pride, anger, impatience, and misplaced zeal damage it, while charity seeks correction at the fitting time. Thus peace is neither the silence of people afraid to speak nor cooperation in something wrong. It is common life under a good that all are called to receive.\footnote{Thomas, \textit{Super Ephesios}, ch.~4, lect.~1, Dessain vol.~2, pp.~312--315.}

\subsection{One people under one Lord}
Thomas compares the Church's unity to that of a city: one ruler, one law, common signs, and a common end. The apostle's Lord, faith, baptism, and Father supply the corresponding order. Faith can name what is believed as well as the habit of believing; baptism is a common sacramental sign whose validity does not depend on a minister's personal holiness. A person's confidence in belonging therefore rests on more than admiration for a particular leader. The shared gifts have an objective source beyond the qualities of the people who administer or receive them.\footnote{Thomas, \textit{Super Ephesios}, ch.~4, lect.~2, Dessain vol.~2, pp.~315--317.}

The Gospel identifies that ruler more fully. David's son does not stand only within an earthly genealogy; David himself calls him Lord. Chrysostom's argument leaves the descent intact while opening the hearer to Christ's divine dignity. The Church's one Lord therefore exceeds every teacher, officeholder, and circle of admirers. The community cannot possess him as a badge distinguishing its members from people it despises. His first command requires whole-hearted love of God, and its like requires love of the neighbor. Lordship governs the use of religious knowledge as well as the content of that knowledge.\footnote{Chrysostom, \textit{Matthew}, hom.~71, exposition of Matt.~22:37--46.}

The Gradual strengthens the claim that the people is received. Augustine understands the blessed nation as belonging to the heavenly city, chosen by God's gift rather than constituting itself by a superior claim. The Lord's possession of it means his care for it. Bellarmine directs its faith, hope, and charity toward the final vision of God. Election is therefore a source of dependence and hope, not a mandate for a modern state's self-exaltation. The psalm's creation praise also places every human society beneath the God whose Word establishes the heavens.\footnote{Augustine, \textit{Psalms}, Ps.~32, first exposition, paragraph 12; Bellarmine, \textit{Psalms}, Ps.~32, commentary on vv.~6 and 12.}

The Collect's request for a pure following belongs here. Allegiance to God alone frees common life from the competing claims of vanity, domination, and faction. The opening servant of the Introit can confess God's just judgment without claiming exemption for his own group. The people that asks for mercy stands under the same law of love it announces to others. Its unity is strengthened by truthful confession, because self-defense at another's expense keeps division alive.

\subsection{A common voice that includes the afflicted}
The Alleluia's singular cry can sound small beside Paul's great list of common gifts. Augustine joins them through the relation of Head and members. Christ took our poverty, and the poor and penitent belong to his body. Consequently the voice of one sufferer can be heard within Christ's common voice without losing its actual suffering. The distress of one generation and then another reaches the same Head. He does not have to commit each member's sin in order to bear the member in his body.\footnote{Augustine, \textit{Psalms}, Ps.~101, paragraphs 1--3.}

This account changes the practical meaning of praying together. The confident member cannot decide that the afflicted member's cry is an interruption of the community's praise. The Alleluia itself holds praise and petition together. Nor does the afflicted person have to turn his need into a general statement before he may belong. The one body can speak in the first person because its members' needs are carried within a common life. Solidarity here is deeper than sympathetic observation from outside.

Daniel's Offertory gives that common voice a named object: the sanctuary and the people over whom God's name has been invoked. His prayer is neither merely a private crisis nor an accusation directed at everyone else. Jerome's two accounts of Daniel's confession preserve the difference between personal guilt and humble identification with the guilty. Both oppose the fantasy that a person's holiness would make intercession unnecessary. One can seek mercy for a community without asserting an equal share in every wrong committed within it.\footnote{Jerome, \textit{Daniel}, II, on Dan.~9:17--20; the chant adapts Dan.~9:17--19.}

The Secret and Postcommunion continue this plural prayer. The holy action is asked to free us from offenses; the mysteries are asked to cure our vices. A member cannot receive the common good while treating the other's need for healing as evidence that the other does not belong. Paul's forbearance becomes sacramentally searching: the assembly asks together for the cure that each member needs in a different way.

\subsection{Common truth and different gifts}
The Communion gathers gift-bearers around the Lord. Augustine asks what it means to stand around a center that is common to all. Truth is not someone's private possession, to be used to draw hearers into a faction. Those who recognize its common source offer in humility. Augustine's image does not make all opinions equally true. It concerns the ownership a proud person falsely claims over truth and the people whom truth should serve.\footnote{Augustine, \textit{Psalms}, Ps.~75, paragraph 12.}

The next verse humbles princes. Augustine reads their spirit as pride and turns the demand for rule toward governance of one's own body. Bellarmine retains the concrete subjection of kings' lives to God. Neither local influence nor earthly sovereignty is the final measure of worth. The community's center belongs to the Lord, who judges its most powerful members as surely as its weakest.

Within that common subjection, different ministries remain real. Chrysostom argues that baptism, salvation through faith, the Spirit, and the Father are shared, while distinct gifts are given for the building of the Church. A larger responsibility therefore adds work and accountability; it is not evidence that another member has received a lesser kind of salvation. Envy mistakes a ministry for a private advantage, and arrogance mistakes a gift for self-created superiority. Both misunderstand why the gift was given.\footnote{Chrysostom, \textit{Ephesians}, hom.~11, exposition of Eph.~4:7--12.}

\subsection{The four senses of a gathered people}
\begin{fourSenses}
\item[Literal.] Israel's elected people, the afflicted psalmist, Daniel's community, Paul's addressees, and Jesus' interlocutors have distinct settings. Their words include communal calling, real supplication, creation praise, and concrete commands. The liturgical assembly receives these different voices; their historical speakers are not identical.

\item[Allegorical.] The Church is the body of David's Lord, established through Word and Spirit, praying with its Head, and nourished by his mysteries. The heavenly city supplies the common destination toward which election and calling are ordered.

\item[Moral.] Shared baptism and hope require humble service across differences. Correction serves charity, patience refuses contempt, and knowledge is offered as a common good. Different tasks remain, but neither envy nor self-exaltation may define the relations among those who perform them.

\item[Anagogical.] The one hope tends toward a shared inheritance and perfected peace. Bellarmine's movement from faith, hope, and charity to vision, and Thomas's account of sacramental unity perfected in glory, give the gathered people an end beyond every temporary alliance or ruler. The Postcommunion asks that present healing bear precisely that lasting fruit.\footnote{Bellarmine, \textit{Psalms}, Ps.~32, commentary on v.~12; Thomas, \textit{Summa}, III, q.~79, a.~2, corpus.}
\end{fourSenses}
